\section{Methodology}

\subsection{Protocol and research questions}
This systematic literature review followed the protocol maintained for SLR1.  The protocol defines a nanosatellite-class population, an MBSE intervention, and outcomes concerned with general modeling, EPS-relevant value, and verification before physical development.  The review method is intentionally narrower than broad space-MBSE syntheses: Buruso et al. \cite{Buruso2026MBSESpace} are used as a methodological and positioning comparator, especially for transparent review procedure and concise enabler classification, but their review is not treated as evidence about the papers in this corpus.

The protocol research questions were:
\begin{description}
\item[RQ1.] How can nanosatellite missions and EPS be modeled in a general manner with minimal knowledge about their behavior?
\item[RQ2.] What are the claimed outcomes related to reliability, reproducibility, and integrability derived from EPS modeling in nanosatellite missions?
\item[RQ3.] How have verification procedures been integrated into EPS models prior to physical development?
\end{description}

The population was interpreted as CubeSat and nanosatellite-class missions rather than larger spacecraft or smaller educational platforms such as CanSats.  For RQ1, the focal intervention was system-level or EPS-level MBSE.  For RQ2, outcomes had to be attributable to MBSE artifacts, MBSE activities, or model-enabled process changes rather than to the selected EPS technology alone.  For RQ3, the review looked for requirements, verification, validation, IVV, test, simulation, or analysis artifacts integrated in or around the model before hardware development.

\subsection{Information sources and search strategy}
The database scope was Scopus and IEEE Xplore.  Scopus was searched directly.  IEEE Xplore coverage was incorporated through manual CSV exports that were screened under the same protocol.  Google Scholar exact-title fallback was used only during the later citation-chasing continuation to confirm or recover specific citation-chasing targets after Scopus misses; it was not used as an open-ended third database.

Search strings combined nanosatellite/CubeSat population terms with MBSE terms such as ``model-based systems engineering'', MBSE, SysML, and selected alternative terminology including Capella, Arcadia, and cyber-physical systems where relevant.  The search strategy followed four concept clusters derived from the protocol: (i) general nanosatellite or CubeSat MBSE reference models, templates, behavior models, and mission/system architectures; (ii) EPS, power-subsystem, power-budget, and electrical-architecture modeling; (iii) verification, validation, IVV, test, requirements-in-the-loop, and traceability; and (iv) outcome and integration terms such as reliability, reproducibility, integrability, interoperability, and lifecycle continuity.  The final candidate-building set combined the non-redundant Scopus backbone, targeted exact-title and citation-chasing recoveries, and the retained IEEE Xplore promotions documented in the search log.

Search and screening limits followed the protocol: publication years 2013--2026, English language, conference or journal papers, and either open access or availability through TEC library access used in this review.  The search log was retained as the audit trail for query identifiers, query strings, result totals, recovery searches, and the stopping rationale.  The stopping decision was based on saturation of the Scopus--IEEE evidence base under the protocol and on the observation that later EPS/V\&V recall passes mainly confirmed sparsity rather than uncovering a large hidden literature.

\subsection{Selection process and eligibility criteria}
Records from the retained search files were merged and deduplicated before screening.  The final PRISMA-style flow was: 199 raw merged records, 118 duplicates removed, 81 unique records screened, 37 excluded during screening, 44 reports sought for retrieval, 4 reports not retrieved or inaccessible, and 40 reports extracted.

Screening applied the protocol criteria as follows:
\begin{itemize}
\item include English conference or journal papers from 2013--2026;
\item include CubeSat, nanosatellite, or very close nanosatellite-class system-design studies, and exclude larger generic spacecraft, payload-only, ground-segment-only, software-only, hardware-calibration, education/tutorial, proceedings-volume, and non-space records unless the screening log recorded a specific protocol-based rationale for retention;
\item include studies using MBSE, SysML, model-based architecture, model-based requirements, model-based simulation/analysis, or closely related digital-engineering approaches at system level, mission level, EPS level, or an EPS-adjacent power/electrical-architecture level;
\item for RQ1, prioritize system-level or EPS-level models, reference models, templates, reusable workflows, behavior models, and architecture models over isolated subsystem models;
\item for RQ2, retain claimed outcomes only when the outcome could be linked to MBSE artifacts or model-enabled processes, not merely to EPS component choices;
\item for RQ3, require explicit verification, validation, IVV, test, requirements traceability, simulation, parametric checking, or pre-physical-development evaluation evidence in or around the model.
\end{itemize}

The screening log distinguishes core and supporting studies.  Core papers directly shaped the synthesis because they were central to CubeSat/nanosatellite MBSE, EPS/power modeling, traceability, reusable reference models, or V\&V.  Supporting papers were retained when they provided contextual evidence on toolchains, lifecycle continuity, requirements-in-the-loop, risk, project-management, or related nanosatellite MBSE practices without being as central to the protocol RQs.

\subsection{Data extraction and quality appraisal}
Extraction was performed on the 40 accessible reports listed in the accessible corpus file, comprising 21 core and 19 supporting studies.  The four reports that remained inaccessible at retrieval were retained in corpus-management artifacts but were not included in the extraction synthesis.

A standardized extraction schema was used for every accessible study.  Extracted fields covered bibliographic and review-management metadata; research-question relevance; mission/platform context; lifecycle stage; EPS relevance; model maturity; MBSE objective, scope, artifacts, abstraction level, and reuse mechanism; implementation pathway and tailoring; toolchain, interoperability, and lifecycle-continuity mechanisms; EPS and verification-specific artifacts; claimed outcomes, mechanisms, evidence types, metrics, baselines, and confidence ratings; governance, ROI, deliverables, and project-decision linkage; gap mapping; limitations; direct evidence quotes; lineage notes; and extraction confidence.

Quality appraisal used the protocol checklist rather than an exclusion threshold.  Each checklist item was scored as No = 0, Insufficiently = 0.5, or Yes = 1.  The checklist assessed venue credibility; whether limitations, threats, assumptions, or scope boundaries were discussed; clarity of nanosatellite/CubeSat context and lifecycle stage; clarity of MBSE method, language, toolchain, model artifacts, and system boundary; whether claimed outcomes were explicitly linked to MBSE artifacts or activities; whether evidence was provided for claims; whether requirements, verification, validation, design decisions, or lifecycle gates were traced to model artifacts when relevant; and whether model availability, reusable assets, data exchange, or implementation details were reported.  Quality scores were used to inform sensitivity, confidence, and discussion of evidence strength; they were not used to automatically remove studies after full-text extraction.

No independent inter-rater reliability exercise is documented in the review artifacts.  Consequently, the review does not report agreement statistics.  Instead, extraction consistency was supported by the fixed protocol, a common extraction schema, explicit rationales and direct quotes for coding decisions, extraction-confidence ratings, and cross-checks against the corpus, screening, PDF-status, and scoping-synthesis artifacts.

\subsection{Synthesis strategy}
The synthesis was structured by the three protocol RQs rather than by a general MBSE tool or method taxonomy.  Descriptive statistics were used to summarize corpus characteristics, study tiers, EPS relevance, approach classifications, quality-score distribution, research-question relevance, value-claim evidence, and V\&V linkage categories.  Narrative synthesis then interpreted these counts in relation to the protocol questions.

A concise enabler taxonomy was retained only as a secondary explanatory layer.  The Pure, Enhanced, Extended, and Integrated MBSE/digital-engineering categories, together with lifecycle activity anchors such as requirements management, model development, simulation/analysis, configuration management, and collaboration/communication, were used as sensitizing categories consistent with the broader space-MBSE review style of Buruso et al. \cite{Buruso2026MBSESpace}.  The main synthesis, however, focused on two issues specific to this review: the evidence behind claimed value and the linkage between MBSE artifacts and V\&V before physical development.

For RQ1, studies were compared by modeling focus: whole-nanosatellite models that include EPS or power as a subsystem, generic whole-nanosatellite models without explicit EPS, and EPS-only or EPS-first models.  Reuse mechanisms such as reference models, templates, profiles, libraries, open packages, and repeatable workflows were examined as evidence for generality and for modeling with limited prior behavioral knowledge.  For RQ2, each value claim was synthesized by outcome domain, specificity, MBSE mechanism, evidence type, metric or value indicator, baseline or comparator, and evidence-confidence rating.  This prevented asserted benefits such as reliability, reproducibility, integrability, schedule reduction, or communication improvement from being treated as measured effects unless the extraction recorded supporting evidence.  For RQ3, V\&V linkage was synthesized across requirements-to-model traceability, verification matrices, parametric and power-budget checks, executable simulation or scenario analysis, test cases and procedures, automated tests, external tool/data-flow coupling, and evidence of physical, software-in-the-loop, hardware-in-the-loop, or integration-test readiness.

\subsection{Review validity}
Several measures were used to strengthen review validity while preserving the limits of the available artifacts.  Construct validity was supported by deriving eligibility, extraction fields, and synthesis categories from the protocol RQs and by keeping the broad space-MBSE review \cite{Buruso2026MBSESpace} as a positioning reference rather than as corpus evidence.  Search validity was supported by using two major engineering bibliographic sources, exact query logging, deduplication, targeted recall searches for EPS/V\&V sparsity, IEEE parity checks, and citation-chasing recovery for specific missing titles.  The main search limitation is that the database scope was restricted to Scopus and IEEE Xplore, with Google Scholar used only for exact-title fallback, so grey literature and some publisher-specific records may be underrepresented.

Selection validity was supported by explicit inclusion/exclusion criteria and a maintained screening log.  Residual uncertainty remains for records with limited abstracts, documented borderline/contextual inclusions, and inaccessible reports.  Extraction validity was supported by full-text extraction for accessible studies, direct evidence quotes, quality appraisal, claim-evidence coding, and extraction-confidence ratings.  The absence of a documented inter-rater exercise is a limitation; no reliability statistic is claimed.

Synthesis validity was addressed by separating descriptive corpus patterns from evidence-backed findings.  Repeated publications from the same CubeSat reference-model lineage were treated with lineage notes so that publication count would not be mistaken for independent replication.  Claims about MBSE value were interpreted through evidence confidence and baseline availability, and EPS-specific conclusions were kept conservative because the corpus contains many system-level or EPS-adjacent studies but few direct EPS-first studies.  External validity is therefore limited to nanosatellite/CubeSat-class MBSE under the search scope and should not be generalized to all space MBSE practice or to all spacecraft EPS development without further evidence.
